\subsection{What is the base space doing?}\label{sec:forget}

A Daniell-style integral is nothing more than a positive linear functional,
satisfying a certain continuity, on a Riesz space of functions on a base space.
From this appearance one might think that the base space is not essential and
could be removed.
In this subsection we see that the role the base space plays is exactly one
thing: \textbf{to supply a place where infima are computed outside $L$}.
We then explain that the direction of ``forgetting the base space'' does not
work as it stands.

\subsubsection{The base space appears in only one place}

In the definition of an integration system $(X,L,I)$, the only place where the
base space $X$ is essentially at work is the hypothesis of the continuity axiom,
that is, the left-hand side of
\[
  f_n\downarrow0\ \text{(pointwise)}\implies I(f_n)\to0 .
\]
If this is replaced by decrease in the order of $L$ itself, then even the
ordinary integral ceases to be an integral.

\begin{example}\label{ex:fat-cantor}
  Let $L$ be the set of all step functions on the closed interval $[0,1]$, that
  is, the set of all functions for which there is a finite partition
  $0=x_0<x_1<\cdots<x_m=1$ such that the function is constant on each
  $[x_{i-1},x_i)$ with $i<m$ and on $[x_{m-1},1]$.
  The partition points are strictly increasing, and singletons are not allowed.
  Taking $I$ to be the ordinary Riemann integral, $(L,I)$ is an integration
  system on $[0,1]$.

  Let $A\subset[0,1]$ be a fat Cantor set of measure $1/2$ and write
  $[0,1]\setminus A=\bigsqcup_kI_k$ with $I_k=(a_k,b_k)$.
  Putting
  \[
    f_n:=1_{[0,1]}-\sum_{k=1}^n1_{[a_k,b_k)}\in L ,
  \]
  the sequence $\{f_n\}$ is monotonically decreasing and $I(f_n)\to1/2$.

  On the other hand, if $g\in L$ satisfies $g\le f_n$ for every $n$ then
  $g\le1_{A\setminus\{a_k\mid k\ge1\}}$ pointwise.
  The set $[0,1]\setminus A$ is a dense open set, and each interval of
  constancy of $g$ is nondegenerate, so it meets some $I_k$.
  At a point of that intersection the upper bound is $0$, so the value of $g$ on
  each interval of constancy is at most $0$, that is, $g\le0$.
  Conversely, since $f_n\ge0$, if $g\le0$ then $g\le f_n$.
  Hence the set of lower bounds of $\{f_n\}$ is $\{g\in L\mid g\le0\}$, whose
  greatest element is $0$.

  That is, in the order of $L$ we have $f_n\downarrow0$, yet
  $I(f_n)\downarrow1/2$.
  In terms of pointwise convergence, $f_n\to1_A\ne0$, and $f_n$ does not
  decrease to $0$.
\end{example}

\noindent
What this example shows is \textbf{the difference between taking the infimum
inside $L$ and taking it outside}.
Taken inside $L$ we get $\inf_nf_n=0$, but taken inside
$\mathbb{R}^{[0,1]}$ we get $\inf_nf_n=1_A$.
What Daniell's continuity axiom requires is the latter, and the base space $X$
supplies $\mathbb{R}^X$, a container large enough for that purpose.

This difference is also at work in the definition of operator morphisms in
Section~\ref{chap:operator}
(Remark~\ref{rem:order-continuity-redundant}).
There too, it is exactly this that makes the statement fail if the order
continuity condition is read in the order of $M$.

\subsubsection{The direction of forgetting the base space does not work as it stands}

One might also consider removing the base space.
In what follows, an \textbf{abstract Riesz space} means an ordered vector space
that is a lattice (Remark~\ref{rem:riesz-terminology}).
The idea is to recover the base space from an abstract Riesz space $L$ by
\[
  \operatorname{Spec}(L):=\operatorname{Hom}_{\mathbf{Riesz}}(L,\mathbb{R}) .
\]
Assigning to each $f\in L$ the function
$\widehat{f}:\operatorname{Spec}(L)\to\mathbb{R};x\mapsto x(f)$ yields a left
adjoint of the functor forgetting the base space.
The adjunction itself does hold. What fails is rather that the unit
$L\to\widehat{L}$ be an isomorphism, that is, that the recovery can actually be
carried out.

There is also the problem that $\operatorname{Spec}(L)$ is too large, but the
fatal point is the following: \textbf{$\operatorname{Spec}$ collapses on exactly
the objects that matter.}

\begin{proposition}\label{prop:spec-collapses}
  Let $L$ be the abstract Riesz space obtained from the set of all Lebesgue
  integrable functions on $[0,1]$ by dividing by the null functions (that is,
  $L^1[0,1]$). Then
  \[
    \operatorname{Hom}_{\mathbf{Riesz}}(L,\mathbb{R})=\{0\} .
  \]
\end{proposition}
\begin{proof}
  Let $\varphi$ be a Riesz homomorphism.
  If $f\ge0$ then $f=f\vee0$ gives
  $\varphi(f)=\varphi(f)\vee0\ge0$, so $\varphi$ is positive.
  A positive linear functional on a Banach lattice is norm bounded
  (\cite{AbramovichAliprantis}), so from
  $L^1[0,1]^*=L^\infty[0,1]$ there is a $g\in L^\infty$ with $g\ge0$ such that
  $\varphi(f)=\int_0^1fg$.

  Suppose $\varphi\ne0$. Then one can take a set $S$ of positive measure on
  which $g>0$.
  Split $S$ into two sets $S_1,S_2$ of positive measure and put
  $f:=1_{S_1}$ and $h:=1_{S_2}$.
  Their supports are disjoint, so $f\wedge h=0$, and since $\varphi$ is a Riesz
  homomorphism,
  \[
    0=\varphi(f\wedge h)=\varphi(f)\wedge\varphi(h)=\min\{\varphi(f),\varphi(h)\} .
  \]
  But $\varphi(f)=\int_{S_1}g>0$ and $\varphi(h)=\int_{S_2}g>0$, a
  contradiction.
  Hence $\varphi=0$.
\end{proof}

\noindent
So $\operatorname{Spec}$ is not ``too large''; it
\textbf{vanishes the moment one divides by the null functions}.
And if one is to do measure theory seriously, dividing by the null functions
cannot be avoided.

Note that treating order continuous functionals on abstract Riesz spaces and
their extensions is itself a classical subject, and anyone pursuing the
direction of this subsection will find \cite{LuxemburgZaanen} and
\cite{FremlinTRS} to be the places to look.

\subsubsection{The correct form, and further prospects}

\begin{remark}[the correct form]
  The idea of ``forgetting the base space and recovering it when needed'' has a
  classical answer.
  For a Daniell closed integration system, the quotient $\widetilde{L}$ of
  Subsection~\ref{sec:ae} is a Banach lattice whose norm is additive on the
  positive cone, that is, an \textbf{abstract Lebesgue space} (AL-space).
  And by Kakutani's representation theorem \cite{Kakutani} (see also
  \cite{Yosida} for the representation of Riesz spaces as function spaces),
  every AL-space is isomorphic to some $L^1(\mu)$.
  That is, the recovery of the base space is already complete here.

  As stated in Subsection~\ref{sec:prior-work} of
  Section~\ref{chap:intro}, however, the more one moves in this direction the
  closer the main theorem of this paper comes to a restatement of an abstract
  completion.
  \textbf{Abstraction and concrete realization stand in a trade-off.}

  The category $\mathbf{Int}_{\mathrm{op},\mathrm{ae}}$ of
  Section~\ref{chap:operator} has in fact already arrived at the end of this
  direction.
  That category distinguishes objects only up to the pre-AL space
  $(\widetilde{L},\widetilde{I})$
  (Subsection~\ref{sec:what-op-ae-sees}), and what the Daniell extension gives
  is its Banach completion.
  That is, both ends of the trade-off sit side by side inside this paper.
\end{remark}

We mention two directions of generalization.

\begin{itemize}
  \item \textbf{Changing the object of values.}
    If $f_n$ is taken to be Banach-space valued while the dominating functions
    $\varphi_n$ remain real-valued, then the proofs of this paper go through
    almost unchanged and one obtains the Bochner integral.
    This is the very subject of \cite{Mikusinski} and is known, but it is a
    technical bridge towards stochastic integration.
  \item \textbf{Internalizing the definition itself.}
    This is the direction of taking the Riesz space and the reals inside a topos
    rather than in the category of sets.
    The ``base space'' is then replaced by the topos itself.
    There is earlier work in this direction.
    \cite{Borger2010} treats measures and integrals inside a monoidal closed
    category and defines the space of simple functions by a universal property.
    The same author's \cite{BorgerHabil} is broader in scope, but we have not
    obtained its text.
    That both are theorems different from this paper's was stated in
    Subsection~\ref{sec:prior-work}.
    There is also \cite{CoquandSpitters} as a constructive, point-free theory of
    integration.
    For presheaf toposes, and especially for presheaves indexed by a filtration,
    a relation to martingales is to be expected, since conditional expectation
    has the shape of an operator morphism of
    Section~\ref{chap:operator} (see also \cite{Simpson}).
\end{itemize}

\noindent
In either direction, the main theorem of this paper retains its meaning only in
so far as it is ``accompanied by a concrete realization''.
The question of how far to abstract, and at what point concreteness is given up,
is one that the framework of this paper raises and leaves open.
